\documentclass[12pt]{article}
\usepackage[utf8]{inputenc}
\usepackage{geometry}
\usepackage{titlesec}
\usepackage{hyperref}
\usepackage{enumitem}
\usepackage{amsmath}
\usepackage{fancyhdr}
\usepackage{cite}

% Page Setup
\geometry{a4paper, margin=1in}
\pagestyle{fancy}
\fancyhf{}
\rhead{Malik Henry}
\lhead{Master's Project Proposal}
\cfoot{\thepage}

% Title Information
\title{\textbf{A Flying Wing Design and Optimization Software Package using AeroSandbox} \\ \large Master's Project Proposal}
\author{Malik Henry \\ mhenry2@aggies.ncat.edu}
\date{November 26, 2025}

\begin{document}

\maketitle

\begin{abstract}
This proposal outlines the development of a flying wing design and optimization software package. The tool will leverage the open-source AeroSandbox Python library, augmented with a custom structural and geometry pipeline to supplement areas where AeroSandbox currently lacks functionality. The primary objective is to create a powerful, user-friendly tool capable of generating and analyzing efficient flying wing designs across various flight regimes, from small-scale sport aircraft to larger heavy-lift platforms.
\end{abstract}

\tableofcontents
\newpage

\section{Introduction and Motivation}
The design of aerodynamically efficient aircraft is a complex, multidisciplinary challenge. Flying wing aircraft, in particular, offer significant potential for high aerodynamic efficiency due to the elimination of parasitic drag from a conventional fuselage and tail. However, their design is notoriously difficult to do properly, presenting unique challenges in stability, control, and structural integration.

Currently, aircraft design is split into two regimes, the preliminary and conceptual design stage with software packages like OpenVSP and XFLR5 but don't have direct pathways into creating parametric feature based geometry. This is in Contrast to tools like SolidWorks and Fusion 360 that can make detailed CAD geometry but lack the fast iteration capabilities of lower fidelity software. This creates a bottleneck in the design process that can lead to substantial lost time if a key aspect of the design needs to be altered due to a change in mission constraints. This leaves a gap in the market for an integrated software tool specifically tailored for the design and analysis of flying wing aircraft. This project aims to fill that gap, it should be noted that this gap applies to other aircraft configurations as well, but I am limiting the scope to the configuration I am most passionate about.

\section{Problem Statement and Target Audience}
The central problem is the lack of a streamlined, open-source workflow for the rapid design iteration of flying wing aircraft. This project will develop a software package to solve this problem targeted towards the following groups:
\begin{itemize}
    \item \textbf{Aerospace Students:} Serving as an accessible platform for senior design and graduate research through rapid, iterative design and immediate feedback.
    \item \textbf{Small to Medium-Sized Enterprises (SMEs) \& Startups:} Enabling rapid prototyping and conceptual design of Unmanned Aerial Vehicles (UAVs) with limited resources.
    \item \textbf{Advanced Hobbyists:} Enabling the community to push the boundaries of range and efficiency by validating complex stability parameters in software, thereby bypassing the risks inherent in experimental tailless configurations and the costs associated with design failure.
\end{itemize}

\pagebreak

\section{Literature Review}

\subsection{Legacy Aerodynamic Analysis Tools} For decades, the conceptual design phase has been anchored by tools developed primarily within academic and governmental institutions. These tools, such as Athena Vortex Lattice (AVL), XFLR5, and OpenVSP, operate largely on potential flow theory. These tools offer fast execution times, but at the cost of fidelity particularly in non-linear regimes \cite{Szabo2022}.

\subsubsection{Athena Vortex Lattice (AVL)} Developed at MIT in 1988, AVL remains a standard for the aerodynamic and flight-dynamic analysis of rigid aircraft. It employs an extended Vortex Lattice Method (VLM) to model lifting surfaces as a grid of horseshoe vortices. AVL is particularly adept at calculating total forces, stability derivatives, and performing eigenmode analysis for dynamic stability \cite{Drela1988}. However, uses inviscid, irrotational flow assumptions that limits its validity to the linear portion of the lift curve. It cannot inherently model skin friction, separation, or stall onset. Furthermore, AVL lacks a native Application Programming Interface (API), which makes integration into external program possible but tedious and less computationally efficient. \cite{Szabo2022}.

\subsubsection{XFLR5} XFLR5 represented an evolution in usability, bridging 2D airfoil analysis with 3D wing design. It integrates the XFOIL code for viscous analysis with a 3D VLM and Panel Method. Unlike pure VLM codes, XFLR5 attempts to account for viscous effects by interpolating 2D airfoil polars onto the 3D wing span, allowing for a better approximation of parasite drag and non-linear lift near stall \cite{Lafifi2020}. Despite its popularity for low-Reynolds-number vehicles, XFLR5 is designed as an interactive GUI-driven tool. The lack of a scriptable backend limits its utility when interfacing with outside programs, Some of these problems have been alleviated in it's successor Flow5, however at time of writing it remains closed source. \cite{Lafifi2020}.

\subsubsection{OpenVSP and VSPAERO} Addressing the geometry bottleneck, NASA’s OpenVSP defines geometry parametrically (e.g., span, chord, sweep) rather than through constructive solid modeling, ensuring shapes are suitable for aerodynamic meshing. It includes VSPAERO, an integrated solver offering both VLM and higher-order Panel Methods, the latter of which allows for the analysis of non-lifting bodies such as fuselages \cite{VSPAERO2024}. A distinct advantage of OpenVSP is its API, which permits external scripts to modify geometry and execute solvers. However, for purely rapid polar generation, studies suggest legacy tools like XFLR5 may still offer speed advantages due to lower overhead \cite{VSPAERO2024}.


\subsection{Flying Wing Aerodynamics and Stability}
The removal of the empennage in flying wing configurations eliminates significant drag but introduces profound challenges in longitudinal and directional stability. To achieve longitudinal trim ($C_{m0} > 0$) without a tail, designers must rely on wing geometry modifications. The most common method involves reflexed airfoils, which curve the trailing edge upward to generate a nose-up moment \cite{MDPI2025}. However, literature indicates that this incurs significant penalties, including reduced maximum lift ($C_{L,max}$) and separation-induced pressure drag, potentially reducing range by over 11\%.

An alternative approach involves geometric twist (washout) and sweep to create a bell-shaped lift distribution, a concept revitalized by the Prandtl-D research \cite{Bowers2023}. Unlike the elliptical distribution which minimizes induced drag for a fixed span, the bell-shaped distribution minimizes induced drag for a fixed root bending moment. Crucially, this distribution sheds vorticity in a way that generates proverse yaw, potentially eliminating the need for vertical tails or inefficient split-drag rudders. Understanding these trade-offs—between the structural weight of a bell-shaped wing and the drag penalties of reflex—is critical for the optimization logic this project seeks to implement.

\subsection{Multidisciplinary Design Optimization (MDAO)}
Flying wings exhibit "strong coupling" between disciplines; the centerbody acts simultaneously as the lift generator, payload bay, and pressure vessel \cite{Harris2013}. Sequential design methods are mathematically incapable of resolving these trade-offs optimally. Therefore, Multidisciplinary Design Optimization (MDAO) frameworks are essential.

Literature distinguishes between monolithic architectures (like Simultaneous Analysis and Design, SAND) and distributed architectures. Monolithic approaches, particularly those leveraging analytic derivatives, have been shown to outperform distributed methods in navigating the narrow feasible design spaces of Blended Wing Bodies (BWB). Furthermore, accurate BWB optimization requires addressing aero-structural coupling. While high-fidelity FEM is often too slow for conceptual sizing, "grey-box" approaches using differentiable beam theory offer a necessary compromise, allowing the optimizer to trade off spar weight against aerodynamic thickness. This project will adopt a monolithic optimization strategy to effectively capture these couplings.

\subsection{The AeroSandbox Library}
The core of this project's software pipeline is AeroSandbox, a modern Python library for automated aircraft design and optimization \cite{Sharpe2024}. It represents a paradigm shift by combining the ease of NumPy with the power of Automatic Differentiation (AD). Unlike finite differencing, AD computes exact gradients analytically, enabling "all-at-once" optimization of aerodynamics, structures, and trajectories at a computational cost comparable to a single function evaluation.

While AeroSandbox excels at sizing and gradient-based optimization, its current capabilities in detailed structural analysis (specifically for complex flying wing internal geometries) and parametric geometry generation are generalized. This project will directly address these limitations by developing a custom geometry and structural pipeline that feeds into the AeroSandbox solvers, bridging the gap between parametric sketching and differentiable physics.

\section{Proposed Methodology}
This project will be executed by developing a Python-based software package. The workflow will consist of three main modules:
\begin{enumerate}
    \item \textbf{Geometry Engine:} A custom module will be developed to procedurally generate parametric flying wing geometries. This will allow for easy manipulation of key design variables such as sweep, twist, and airfoil shape.
    \item \textbf{Analysis Pipeline:} This module will integrate the AeroSandbox library for aerodynamic analysis with a custom-built structural model (e.g., using beam theory or a finite element framework) to estimate structural mass and deflection.
    \item \textbf{Optimization Wrapper:} An outer loop will be created to run MDAO studies. This will allow a user to define objective functions (e.g., maximize lift-to-drag ratio, minimize weight) and constraints (e.g., maintain static stability, do not exceed maximum stress) to find an optimal design.
\end{enumerate}

\section{Actionable Deliverables}
The successful completion of this master's project will yield the following deliverables:

\textbf{Deliverable 1: Functional Software Package} \\
A well-documented, open-source Python program capable of performing MDAO for flying wing aircraft. The software will be hosted on a public repository (e.g., GitHub) and will include a user guide with example cases.

\textbf{Deliverable 2: Validated Aircraft Designs} \\
The software will be used to generate a complete aircraft designs, demonstrating its operability:
\begin{itemize}
    \item A 10-foot wingspan heavy-lift aircraft, optimized for maximum payload capacity, Design constriants will be based upon SAE Aero Design 2026 Regular Class design rules
\end{itemize}
Complete design files and performance predictions for the aircraft will be provided.

\textbf{Deliverable 3: Computational and Experimental Validation} \\
The performance predictions for the two designs will be validated through:
\begin{itemize}
    \item \textbf{CFD Analysis:} High-fidelity Computational Fluid Dynamics (CFD) simulations will be run on the final designs to verify the aerodynamic performance predicted by the software.
    \item \textbf{Real-World Flight Tests:} The smaller sport aircraft will be constructed and flight-tested to provide qualitative and quantitative validation of its handling characteristics and performance.
\end{itemize}
\pagebreak
\section{Summary}
This project proposes the development of an open-source, Python-based software package for the design and optimization of flying wing aircraft. By integrating the AeroSandbox library with custom geometry and structural modules, the tool will enable rapid, multidisciplinary design iteration—addressing a critical gap in current conceptual design workflows. Targeted at aerospace students, SMEs, startups, and advanced hobbyists, the software will support the generation of efficient flying wing configurations through gradient-based MDAO. Deliverables include a fully documented codebase, a validated aircraft design, and both computational (CFD) and experimental (flight test) validation. The outcome will be a practical, accessible platform that bridges low-fidelity conceptual tools and high-fidelity CAD, accelerating the development of flying wing aircraft. 
\newpage
\begin{thebibliography}{9}

\bibitem{Harris2013}
Harris, C., et al. (2013). \textit{Multidisciplinary Unmanned Combat Air Vehicle (UCAV) System Design Using Multi-Fidelity Models}. ResearchGate. https://www.researchgate.net/publication/257511056

\bibitem{Sharpe2021}
Sharpe, P. D. (2021). \textit{AeroSandbox: A Differentiable Framework for Aircraft Design Optimization} (Master's thesis). Massachusetts Institute of Technology. https://dspace.mit.edu/handle/1721.1/140023

\bibitem{Bowers2023}
Bowers, A. H., et al. (2023). \textit{Wake-Structure of the Prandtl-D}. Aerospace Research Central. https://arc.aiaa.org/doi/pdfplus/10.2514/6.2023-4003

\bibitem{Lafifi2020}
Lafifi, M. M. (2020). \textit{XFLR5 Analysis of foils and wings operating at low Reynolds numbers}. ResearchGate. https://www.researchgate.net/publication/348716216

\bibitem{Sharpe2024}
Sharpe, P. D. (2024). \textit{Accelerating Practical Engineering Design Optimization with Computational Graph Transformations} (Ph.D. thesis). Massachusetts Institute of Technology. https://dspace.mit.edu/handle/1721.1/157809

\bibitem{MDPI2025}
Zhang, Y., et al. (2025). \textit{Aerodynamic Design Optimization for Flying Wing Gliders Based on the Combination of Artificial Neural Networks and Genetic Algorithms}. MDPI. https://www.mdpi.com/2226-4310/12/9/818

\bibitem{SUAVE}
Botero, E., et al. (2025). \textit{Current Capabilities and Challenges of NDARC and SUAVE for eVTOL Aircraft Design and Analysis}. Stanford University. https://suave.stanford.edu/publications

\bibitem{Szabo2022}
Szabo, A., et al. (2022). \textit{Comparing Potential Flow Solvers for Aerodynamic Characteristics Estimation of the T-Flex UAV}. International Council of the Aeronautical Sciences. https://www.icas.org/icas\_archive/ICAS2022/data/papers/ICAS2022\_0367\_paper.pdf

\bibitem{Drela1988}
Drela, M., \& Youngren, H. (1988). \textit{AVL User Primer}. MIT AeroAstro. (Contextual citation based on standard AVL literature referenced in source text).

\bibitem{VSPAERO2024}
TechScience. (2024). \textit{Validation of VSPAERO for Basic Wing Simulation}. RIMNI. https://cdn.techscience.cn/files/rimni/2024/online/RIMNI1012/TSP\_RIMNI\_56492/TSP\_RIMNI\_56492.pdf

\end{thebibliography}

\end{document}